% ============================================================
% empirical_strategy.tex
% Sección: Estrategia Empírica
% Tesis: Impacto del Gasto en Limpieza Pública sobre
%        Gestión de Residuos Sólidos Municipales, Perú 2015-2019
% Autor: Dante Barreto Gamarra (UNAS)
% ============================================================

\section{Estrategia Empírica}
\label{sec:estrategia}

\subsection{Identificación y Modelo Base}
\label{sec:identificacion}

La estrategia empírica principal utiliza un diseño de Efectos Fijos de Doble Vía (\textit{Two-Way Fixed Effects}, TWFE) con Correlated Random Effects (CRE) y una función de control (CF) para corregir la endogeneidad del gasto municipal (\citealt{Wooldridge2015}).
El modelo general para el municipio $i$ en el año $t$ es:

\begin{equation}
  Y_{it} = \alpha + \beta \ln G_{it} + \gamma \bar{G}_{i} + \delta_{t} + \varepsilon_{it}
  \label{eq:modelo_base}
\end{equation}

\noindent donde $Y_{it}$ es el indicador de gestión de residuos sólidos (FR, QPRS, CSRS, PI, RS o PRS),
$\ln G_{it} = \ln(\text{Gasto}_{it} + 1)$ es el logaritmo del gasto devengado en limpieza pública,
$\bar{G}_{i}$ es la media temporal de $\ln G_{it}$ para el municipio $i$ (\textit{Mundlak device}),
y $\delta_{t}$ son efectos fijos de año.

\subsection{Corrección de Endogeneidad: CRE + Función de Control}
\label{sec:cre_cf}

El gasto municipal puede estar correlacionado con factores no observables (capacidad institucional, preferencias por el medio ambiente), lo que genera endogeneidad.
Siguiendo \citet{Wooldridge2015}, implementamos la corrección en dos etapas:

\textbf{Etapa 1} (OLS agrupado con Mundlak means):
\begin{equation}
  \ln G_{it} = \pi_0 + \pi_1 \bar{G}_{i} + \sum_t \pi_t \delta_t + v_{it}
  \label{eq:stage1}
\end{equation}

\noindent Los residuos $\hat{v}_{it}$ capturan la variación del gasto no explicada por la heterogeneidad municipal permanente.

\textbf{Etapa 2}: Para los modelos no lineales (Ordered Probit, Probit, Fractional Logit), $\hat{v}_{it}$ se incluye como regresor adicional.
Si $\hat{v}_{it}$ es estadísticamente significativo, confirma la presencia de endogeneidad y justifica la corrección CF.

Los errores estándar se agrupan a nivel provincial (\textit{province-clustered SE}) para corregir la correlación serial y la heteroscedasticidad a nivel de municipio.

\subsection{Supuestos de Identificación}
\label{sec:supuestos}

\subsubsection*{Stable Unit Treatment Value Assumption (SUTVA)}

Para garantizar la validez de las estimaciones CRE+CF, se asume el cumplimiento del supuesto SUTVA.
Dado que la gestión de residuos y la ejecución presupuestal en el Perú son de competencia estrictamente distrital según la Ley Orgánica de Municipalidades (Ley N° 27972), se asume que el gasto del distrito $i$ no genera efectos de desbordamiento (\textit{spillovers}) que afecten los resultados de gestión del distrito $j$.
Esta premisa es razonable porque: (i) los servicios de limpieza pública son no exportables por definición jurisdiccional; (ii) la infraestructura de disposición final (rellenos sanitarios) es mayoritariamente de escala distrital en el período 2015-2019; y (iii) los contratos de concesión de residuos no cruzaban límites distritales en el período de análisis según los registros RENAMU.

\subsubsection*{Exogeneidad del Instrumento FONCOMUN}

Como verificación de robustez, se utiliza la transferencia FONCOMUN como instrumento para $\ln G_{it}$.
La restricción de exclusión requiere que FONCOMUN afecte la gestión de residuos únicamente a través del gasto en limpieza pública, y no directamente.
Siguiendo \citet{ConleyHansenRossi2012}, se reportan límites de \textit{plausible exogeneity} para evaluar la robustez de la restricción de exclusión a desviaciones locales.

\subsubsection*{Tendencias Paralelas (TWFE)}

La validez del TWFE requiere que, en ausencia de variación en el gasto, los municipios habrían seguido trayectorias paralelas en las variables de resultado.
Se verifica este supuesto mediante pruebas de pre-tendencias usando el rezago $\ln G_{it-1}$ (especificación DL(1)) y el análisis de event-study presentado en el Apéndice~\cref{sec:appendix_pretendencias}.

\subsection{Modelos por Dimensión}
\label{sec:modelos_dimension}

\subsubsection*{D1 --- Eficiencia en Recolección}

Para la frecuencia de recolección (FR: 1=diaria, 5=nunca), estimamos un Ordered Probit CRE+CF dado el carácter ordenado de la variable.
Para la cantidad recolectada (QPRS, en kg/día), utilizamos un modelo OLS log-log con indicador de recolección cero (\textit{corner solution}).
Para la cobertura del servicio (CSRS: 1=\textless{}25\%, 4=\textgreater{}75\%), estimamos un Ordered Probit CRE+CF análogo al de FR.

\subsubsection*{D2 --- Planeamiento Municipal}

El índice de planeamiento (PI: suma de 5 instrumentos binarios, rango 0-5) se estima con OLS CRE+CF.
Adicionalmente, cada componente binario (PIGARS, PMRS, SRRS, PTRS, PSFRSRS) se modela con Probit CRE+CF para identificar qué instrumentos responden al gasto.

\subsubsection*{D3 --- Disposición Final Adecuada}

Para el uso de relleno sanitario (RS: binario), se usa Probit CRE+CF.
Para la proporción de residuos enviados a relleno sanitario (PRS $\in [0,1]$), se aplica un Fractional Logit \citep{PapkeWooldridge1996} con la función de control integrada.

\subsection{Corrección por Multiplicidad}
\label{sec:bonferroni}

Las tres hipótesis primarias (H1: efecto sobre FR; H2: efecto sobre PI; H3: efecto sobre PRS) se contrastan simultáneamente con la corrección Bonferroni-Holm para controlar la tasa de error familiar (\textit{familywise error rate}).
Las comparaciones secundarias dentro de cada dimensión (componentes D2, submodelos D1) se reportan sin corrección de multiplicidad, siguiendo la práctica estándar para análisis secundarios preregistrados.

\subsection{Inferencia con Bootstrap de Dos Etapas}
\label{sec:bootstrap}

Para los modelos no lineales de segunda etapa (Ordered Probit, Probit y Fractional Logit), los errores estándar analíticos de la regresión subestiman la variabilidad muestral real porque no incorporan el error de estimación de $\hat{v}_{it}$ generado en la primera etapa \citep{Wooldridge2015}.
Para corregir este sesgo de inferencia, se implementa un bootstrap de dos etapas con $B = 999$ réplicas, remuestreando por conglomerado provincial en ambas etapas simultáneamente: en cada réplica se re-estima la primera etapa (OLS con medias Mundlak) y se obtiene un nuevo $\hat{v}_{it}^{(b)}$, el cual se incorpora en la segunda etapa correspondiente.
Los errores estándar reportados en la \cref{tab:resultados_principales} y en las tablas de componentes son los desvíos estándar de la distribución bootstrap de los coeficientes y efectos marginales promedio (AME).
Para los modelos OLS (QPRS y PI), los errores estándar analíticos clustered a nivel provincial son asintóticamente válidos \citep{Wooldridge2015} y se reportan directamente; el bootstrap confirma la concordancia con los errores analíticos en ambos casos.

\subsection{Frontera de Identificación}
\label{sec:frontera}

La variación identificante en el estimador CRE+CF proviene de las fluctuaciones dentro-municipio del gasto, una vez descontada la media temporal del municipio ($\bar{G}_{i}$): formalmente, el estimador explota $\tilde{G}_{it} = \ln G_{it} - \bar{G}_{i}$, es decir, los desvíos transitorios del gasto respecto al nivel histórico de cada municipio.
Este componente está purificado de la heterogeneidad permanente no observable ---capacidad institucional estructural, geografía, historia fiscal--- que el dispositivo de Mundlak absorbe en la primera etapa.

La frontera de identificación que el diseño no puede superar son los confundidores tiempo-variantes correlacionados con $\tilde{G}_{it}$.
El ejemplo arquetípico es un cambio en la capacidad política local: un nuevo alcalde técnicamente competente puede incrementar simultáneamente el gasto en limpieza pública y la calidad de la gestión de residuos en el mismo período, generando una correlación espuria que la función de control no elimina porque el residuo $\hat{v}_{it}$ no distingue entre un shock de gasto exógeno y un shock de capacidad de gestión.
Los resultados robustecedores abordan este riesgo desde distintos ángulos: el test DL(1) descarta tendencias preexistentes y efectos de anticipación; la significatividad de $\hat{v}_{it}$ en todos los modelos confirma la dirección del sesgo; y los límites de \citet{Oster2019} cuantifican la magnitud del sesgo de variable omitida estática necesaria para anular los efectos (Apéndice~\ref{sec:appendix_oster}).
Sin embargo, ninguna de estas pruebas descarta de forma definitiva un confundidor tiempo-variante coincidente.
Los efectos estimados deben interpretarse como efectos condicionales bajo el supuesto CRE ---que no persisten tendencias temporales heterogéneas no observadas entre municipios--- más que como efectos causales locales en sentido estricto.
La integración de la transferencia FONCOMUN como instrumento externo es la extensión identificatoria que permitirá recuperar el LATE para el subgrupo de municipios cuyo gasto responde a variaciones en esa transferencia, transformando los efectos condicionales aquí reportados en efectos causales locales verificables.
